\section{Discussion}
\label{section:discussion}

\subsection{What the Evidence Supports}
\label{subsec:discussion_key_findings}

Read together, the taxonomy (Section~\ref{sec:taxonomy}), the compression review
(Section~\ref{sec:compression}), and the audit of the evaluation record
(Section~\ref{sec:reported}) support four claims.

\textbf{Capability per parameter has risen sharply.} The strongest 3--4B models
of 2026 report scores on mathematics, code, and knowledge benchmarks that were
state of the art for 13B-class systems in 2023. Within families that publish
controlled ablations, the dominant explanatory variable is training data rather
than architecture: corpus composition and curation move results more than any
architectural choice available at this scale. The practical implication for model
builders is that data pipelines dominate architectural novelty as a lever on
\gls{slm} quality.

\textbf{Compression is a first-class capability, not an afterthought.} INT8 is
effectively lossless (\evB) and INT4 costs a small, non-uniform amount of
accuracy, with smaller models and mathematical reasoning at the higher end
(\evC). Because single-stream autoregressive decoding is memory-bandwidth-bound,
4-bit quantization is simultaneously a storage optimization and a decode-time
speedup on edge hardware, in proportion to the bandwidth it frees and subject to
the regime conditions of Section~\ref{subsec:deployment_calc_vs_meas}.
Quantization preserves fidelity better than pruning below 4B parameters
(\evB, \citealp{kurz2025pruning_vs_quant}), which simplifies the practitioner's
decision: quantize aggressively, and prune only when a structured size reduction
is independently required. What the literature does not support, and what
Section~\ref{subsec:compression_reporting} withdraws, is a single accuracy-retention
figure per technique: retention depends on baseline, model, size, task, and
aggregation rule, and a number stated without them is not interpretable.

\textbf{Right-sizing beats scaling for bounded tasks.} A quantized 3.8B model
that clears a task's accuracy threshold, running on hardware that already exists
in the user's hand, is preferable on cost, latency, energy, and data exposure to
a larger model that also clears it. A 3--4B model at INT4 occupies roughly 2~GB
and runs at interactive speed on a phone or a Jetson-class module. The
scaling-era assumption that a larger model is always the safer default is
inverted for any task whose requirements can be stated in advance \emph{and
checked}. Both halves of that condition matter: where the requirement cannot be
stated in advance, the task is not bounded, and where success cannot be checked,
per-call errors compound silently across an agent trajectory
(Section~\ref{subsubsec:agentic_ev_e2e}). Section~\ref{subsec:agentic_counterexamples}
lists the workload classes where this claim does not apply.

\textbf{The evaluation record is the weak point, and it is weaker than the field
acknowledges.} Section~\ref{sec:reported} documents why a cross-model accuracy
comparison cannot be built from current public reporting: benchmark sets do not
intersect across vendors, shot counts vary per benchmark and per vendor, base and
instruction-tuned figures are conflated, and matched baselines are rarely
published. This is not a limitation peculiar to this survey. It applies equally
to the comparison tables circulated in secondary sources, community leaderboards,
and prior surveys, which are widely reproduced without these caveats and which
practitioners use to make procurement decisions. Establishing a shared,
independently run evaluation for this model class would be a more valuable
contribution to the field than another survey, and
Section~\ref{subsec:challenges_verification} specifies what it would require.

\subsection{Implications for Agentic Systems}
\label{subsec:discussion_agentic}

The agentic argument of Section~\ref{sec:agentic} rests on these claims rather
than restating them. Capability-per-parameter growth is what makes an \gls{slm}
\emph{sufficient} for a delegated sub-task; compression and edge deployability are
what make it \emph{cheap} and \emph{private}; right-sizing is the design
discipline that turns a heterogeneous SLM--LLM system from an option into the
default. Two cross-cutting observations sharpen the picture. First, routing must
be capability-aware, not size-aware: the reported multilingual and safety gaps
(Section~\ref{subsec:reported_safety}) mean that a naive ``small model unless
hard'' router will silently degrade non-English and safety-sensitive calls.
Second, adaptation matters more as models get smaller, so the economical unit in
an agent is often a \emph{fine-tuned} small model per recurring call type rather
than a single general \gls{slm}, which is exactly what the LLM-to-SLM conversion
pipeline of Section~\ref{subsec:agentic_conversion} produces.

\subsection{Democratization and Sustainability}
\label{subsec:discussion_implications}

Deploying a 3B model that answers grade-school mathematics correctly on an
\$80 single-board computer is a qualitative change in access. Organizations facing
cloud-API cost, data-sovereignty rules, or unreliable connectivity can now run
capable \gls{nlp} locally with no per-query charge and, subject to the scope
conditions of Section~\ref{subsec:deployment_privacy_scope}, without transferring
prompt content to a processor outside their control.

The environmental case is narrower than it is usually stated and we state it
narrowly. Reported per-inference energy for a small model on modest hardware is
one to two orders of magnitude below a frontier-model query in a datacenter, so
for a task the small model completes correctly, serving it small uses
substantially less inference energy. Whether a wholesale shift reduces total
system energy is a different question, and one no published accounting answers:
it turns on teacher-training energy, the compression pipeline, device embodied
carbon, regional grid intensity, update distribution, and the rebound effect, none
of which the available reporting covers.
Section~\ref{subsubsec:challenges_env_scope} sets out that accounting boundary in
full. We therefore make the per-query claim and decline the aggregate one.

\subsection{Limitations}
\label{subsec:discussion_limitations}

Beyond the evidence-base limitations already stated
(Section~\ref{subsec:reported_limits}), six bound this work.

\emph{No original measurement.} No benchmark was executed for this survey. Every
quantitative claim is therefore at best grade~\evB{}
(Section~\ref{subsec:reported_grading}), and the paper contains no grade~\evA{}
evidence. This is a deliberate choice, since we judged an audit of the record more
useful than adding another unreplicated table to it, but it is a limitation and
Section~\ref{subsec:challenges_verification} is written to be acted on rather than
cited.

\emph{The definitional boundaries are conventions.} The twenty-billion threshold,
the split into conditions C1 and C2, and the sub-10B restriction on the evaluation
cohort are pragmatic choices, not natural kinds. They are stated precisely
(Section~\ref{subsec:intro_definition}) so that a reader can recompute under
different ones, which is the most a threshold of this sort can offer.

\emph{The thesis is a hypothesis.} The agentic argument is grounded in a
production-scale position paper \citep{belcak2025slm_agentic}, two large
deployments, and the agent-specific evidence of
Section~\ref{subsec:agentic_evidence}, and its premises are graded individually in
Table~\ref{tab:agentic_evidence}. The end-to-end claim, that a small-model agent
completes realistic tasks at comparable rate and lower cost per completed task,
is the one premise for which we found no evidence either way. It is a projection,
and we mark it as one.

\emph{Coverage is English-centric.} The task coverage, the benchmarks audited, and
most of the compression literature are English-centric, which matters more here
than usual because compression degrades non-English performance disproportionately
(Section~\ref{subsec:discussion_ethics}).

\emph{Domain sections are uneven.} Code and healthcare have public benchmarks and
are treated evidentially; legal has neither, and
Section~\ref{subsec:applications_legal} states a structural argument rather than a
capability claim. Readers should not read the four verticals as equally supported.

\emph{The snapshot ages.} The ecosystem moves faster than a review cycle. The
update policy of Section~\ref{subsubsec:intro_snapshot} and the versioned deposit
are how we bound this, not eliminate it.

\subsection{Ethical Considerations}
\label{subsec:discussion_ethics}

Safety alignment is proportionally harder at small scale: several research-grade
\glspl{slm} lack the safety training of their larger siblings, and fine-tuning can erode
what alignment exists \citep{qi2023finetuning}. Practitioners deploying SLMs, and
especially those wiring them into agents with tool access, should treat
application-level safety filtering and output monitoring as mandatory, not
optional. Compression carries its own risk: pruning and quantization can degrade
performance unevenly across dialects, languages, and demographic groups
\citep{blodgett2020language}, and the reported advantage of quantization over
pruning is largest precisely on multilingual tasks \citep{kurz2025pruning_vs_quant},
which suggests pruning disproportionately harms non-English users. Fairness-aware
compression objectives and per-group performance floors for high-stakes
applications are open requirements. Against these risks, on-device \glspl{slm} advance
privacy (one egress channel removed, subject to Section~\ref{subsec:deployment_privacy_scope}), accessibility (no connectivity required),
and cost (no subscription), which matters most for the healthcare and
educational applications of Section~\ref{sec:applications} in under-resourced
settings.
